\documentclass[aspectratio=169]{beamer}
\usetheme{Madrid}
\usecolortheme{whale}

\usepackage[utf8]{inputenc}
\usepackage[spanish]{babel}
\usepackage{amsmath,amssymb}
\usepackage{graphicx}
\usepackage{listings}
\usepackage{booktabs}

\lstset{
    language=Python,
    basicstyle=\ttfamily\footnotesize,
    keywordstyle=\color{blue}\bfseries,
    stringstyle=\color{red},
    commentstyle=\color{gray}\itshape,
    breaklines=true,
    showstringspaces=false
}

\title[INT210 - Sesión 0.2]{Sesión 0.2: Estructuras Multicamino y Diccionarios de Despacho}
\subtitle{Módulo 0: Fundamentos Algorítmicos en Python}
\author[Diplomado en Ciencia de Datos]{Cátedra de Inteligencia Artificial y Ciencia de Datos}
\institute[INT210]{Machine Learning From Scratch}
\date{Semana 1}

\begin{document}

\frame{\titlepage}

\begin{frame}{Contenido de la Sesión}
    \tableofcontents
\end{frame}

\section{Introducción y Analogía Postal}

\begin{frame}{La Analogía Infantil: La Central de Clasificación Postal}
    \begin{columns}
        \column{0.55\textwidth}
        \begin{itemize}
            \item \textbf{Enfoque Lineal (\texttt{if-elif}):} El cartero pregunta buzón por buzón de forma secuencial ($\mathcal{O}(N)$).
            \item \textbf{Enfoque Indexado (Despacho Hash):} El código postal dirige la carta instantáneamente por un tubo neumático ($\mathcal{O}(1)$).
        \end{itemize}
        \column{0.45\textwidth}
        \begin{alertblock}{Principio de Escalabilidad}
            A medida que el número de rutas crece, la búsqueda lineal degrada la latencia. La búsqueda indexada mantiene un tiempo de respuesta constante.
        \end{alertblock}
    \end{columns}
\end{frame}

\section{Análisis de Complejidad Algorítmica}

\begin{frame}{Complejidad Temporal: $\mathcal{O}(N)$ vs $\mathcal{O}(1)$}
    \begin{columns}
        \column{0.5\textwidth}
        \textbf{Bifurcaciones Lineales:}
        \begin{equation*}
            T_{\text{lineal}}(N) = \frac{N + 1}{2} \cdot c \implies \mathcal{O}(N)
        \end{equation*}
        \begin{itemize}
            \item Evaluación secuencial de cada condición.
            \item Peor caso: $N$ comprobaciones.
        \end{itemize}
        \column{0.5\textwidth}
        \textbf{Diccionarios Hash:}
        \begin{equation*}
            T_{\text{hash}}(N) = c' \implies \mathcal{O}(1)
        \end{equation*}
        \begin{itemize}
            \item Acceso directo por cálculo hash: $h(\text{key}) \pmod M$.
            \item Independiente del número de elementos.
        \end{itemize}
    \end{columns}
\end{frame}

\section{Structural Pattern Matching: match-case}

\begin{frame}[fragile]{Structural Pattern Matching (Python 3.10+)}
\begin{lstlisting}
match evento:
    case ('MALWARE', severidad) if severidad >= 8:
        return aislar_host()
    case ('PHISHING', _):
        return bloquear_dns()
    case ('DOS', _):
        return aplicar_rate_limit()
    case _:
        return registrar_log_auditoria()
\end{lstlisting}
\begin{exampleblock}{Ventajas}
    Permite desestructurar tuplas, validar tipos y aplicar cláusulas de guarda (\textit{guards}) en una sola construcción sintáctica limpia.
\end{exampleblock}
\end{frame}

\section{Patrón de Despacho en Producción}

\begin{frame}[fragile]{Patrón Dispatch Dictionary en Ciberseguridad y Arte}
\begin{lstlisting}
DESPACHO_ACCIONES = {
    'RANSOMWARE': aislar_host,
    'DDOS': propagar_bgp_flowspec,
    'SQLI': inyectar_regla_waf
}

def procesar_evento(tipo_alerta, datos):
    manejador = DESPACHO_ACCIONES.get(tipo_alerta, accion_por_defecto)
    return manejador(datos)  # Ejecucion O(1)
\end{lstlisting}
\begin{block}{Beneficio de Mantenibilidad}
    Añadir una nueva acción no modifica la función principal (cumplimiento del principio \textit{Open/Closed} de SOLID).
\end{block}
\end{frame}

\begin{frame}{Conclusiones}
    \begin{itemize}
        \item Las cadenas extensas de \texttt{if-elif} deben evitarse en favor de diccionarios de despacho o \texttt{match-case}.
        \item La complejidad $\mathcal{O}(1)$ es esencial para sistemas de seguridad en tiempo real y motores de renderizado.
        \item \textbf{Siguiente Sesión (0.3):} Bucles Iterativos Deterministas (\texttt{for}), \texttt{enumerate} y Comprensiones Pythonicas.
    \end{itemize}
\end{frame}

\end{document}
